% p177_1; source PDF page 177; manually verified editable table.
\begin{tikzpicture}[x=1cm,y=-1cm,inner sep=0pt,outer sep=0pt]
\path[use as bounding box] (0,0) rectangle (16.000,3.050);
\draw[line width=.4pt] (0.000,0.000) rectangle (8.000,0.850);
\node[align=center,text width=7.740cm,font={\linespread{1}\fontsize{10}{13.50}\selectfont\bfseries}] at (4.000,0.425) {法条规定“多次”定一罪的两种情形};
\draw[line width=.4pt] (8.000,0.000) rectangle (16.000,0.850);
\node[align=center,text width=7.740cm,font={\linespread{1}\fontsize{10}{13.50}\selectfont\bfseries}] at (12.000,0.425) {典型规定};
\draw[line width=.4pt] (0.000,0.850) rectangle (8.000,1.950);
\node[align=left,text width=7.740cm,font={\linespread{1}\fontsize{10}{13.50}\selectfont}] at (4.000,1.400) {定一罪，“多次”是犯罪成立条件};
\draw[line width=.4pt] (8.000,0.850) rectangle (16.000,1.950);
\node[align=left,text width=7.740cm,font={\linespread{1}\fontsize{10}{13.50}\selectfont}] at (12.000,1.400) {多次盗窃，多次抢夺，多次敲诈勒索};
\draw[line width=.4pt] (0.000,1.950) rectangle (8.000,3.050);
\node[align=left,text width=7.740cm,font={\linespread{1}\fontsize{10}{13.50}\selectfont}] at (4.000,2.500) {定一罪，“多次”是法定刑升格条件};
\draw[line width=.4pt] (8.000,1.950) rectangle (16.000,3.050);
\node[align=left,text width=7.740cm,font={\linespread{1}\fontsize{10}{13.50}\selectfont}] at (12.000,2.500) {多次抢劫};
\end{tikzpicture}
